% GPT-6 Astra Pro; reasoning effort not exposed. SIGN_NEXT, 2026-11-21.
\subsection{SF-Q6-08: a pointwise inert estimate with the actual input gates}

\paragraph{Status and statement.}
The following is a conditional theorem about the \emph{actual} column
chain and terminal period of an actual datum $D:\texttt{D7}$.
It does not assert that the numerical candidate below supplies all its
hypotheses.  Its purpose is to remove the irrelevant global size and
split-bound hypotheses from one fixed inert-prime estimate.

Fix a word $u$, a coefficient place $w$ of residue characteristic
$p\ge5$, and a used twist $\mu$.  Assume
\texttt{IsResidueClassZero D p}, $h_P(D,p)=0$, and
\texttt{LevelFree Dx u mu w}.  In particular $p$ is distinct from the
column level prime and both type primes, and does not belong to the
exceptional twist set $\texttt{Dx.Qp mu}$.
Use the actual twisted forms $T:\texttt{TwForms D}$ and exact kernel
entry data $E:\texttt{OrbitPlaces.Ent D}$.  Write
$s_{i,z}=|\nu_{i,z}|$ for the sizes of the partitions in $E$, and
$e_w$ for the ramification index normalizing the source's
$\texttt{ur D w}(b)=e_wb$.  The required gates are:
\begin{enumerate}
\item For every step $i<m$ and every actual pair of column classes,
      the series $E.\texttt{Th}$ at this fixed $w$ satisfies
      \texttt{ThetaExpansionRealPlaces}.  This includes the integral
      remainder and the logarithmic theta factorization with conditions
      B1--B4; coefficient algebraicity alone is insufficient.
\item The actual transfer statement \texttt{KForm E Dx T 1} holds at
      this word, twist and place, for all steps and all bounds.  This
      transfers the Schur coefficient estimate to the actual column
      lattices, including class sums and stabilizer denominators.
\item The actual terminal pairing satisfies \texttt{PairInt D Dx T}
      at this word, twist and place, and \texttt{DGamDef D} holds.
      The row vector lies in the specified lattice
      $\widetilde M(\mu)$ and space $V_0(\mu)$.
\item $N_iN_{i+1}$ is even for each step, and a real number $A$ satisfies
      $p\le A$ and $s_{i,z}\le A$ for all $i,z$.
\end{enumerate}
Let $d_w=\texttt{vDG D u w}$, with the source normalization.
For every input $\phi\in\texttt{MramL D u}$ and every indicated row
vector $F$, the actual period is the image of some $y$ in the actual
coefficient field.  Either $y=0$, or
\begin{equation}\label{eq:SF08-pointwise}
 \frac{v_w(y)}{e_w}\ \ge\
 \sum_{i<m}\left(
   \frac{\sum_z s_{i,z}}{p^2-1}
       -fN_iN_{i+1}(\log_p A+1)\right)-m d_w.
\end{equation}
The same conclusion applies to any specified input lattice locally
contained in \texttt{MramL}, using that containment explicitly.
No nonvanishing of the period is asserted.

\begin{proof}
The residue and $h_P=0$ assumptions give $n_P=1$.
For each $z$, the one-orbit local theorem
\texttt{Loc1.ThmOne 1} applies.  Its entire
\texttt{Loc1.HidaSetting A} hypothesis is
$p\le A$ and $|\nu_{i,z}|\le A$.
In particular it does not require a logarithmic lower bound on the
matrix rank.  Its estimate is
\[
 \frac{s_{i,z}}{p^2-1}
       -N_iN_{i+1}(\log_p A+1).
\]
The closed source theorem
\texttt{split\_bound\_eleven\_one} supplies this constant-one local
bound.  The theorem \texttt{orbitsAddAtConjugate\_ent} adds these
bounds over the $f$ real places using the stated theta expansion;
thus it gives \texttt{EntGe} with
\[
 B_i=\frac{\sum_zs_{i,z}}{p^2-1}
                   -fN_iN_{i+1}(\log_p A+1).
\]
The actual \texttt{KForm} gate with $\texttt{sp}=1$ gives
\texttt{KGe} with $B_0$ at the first step and $B_i-d_w$ at
all subsequent steps.  The proof of \texttt{OrbitPlaces.compose}
iterates these actual lattice estimates.  Its twisted branch uses
\texttt{base\_map\_mem\_local\_lattice}; consequently the twist
and its integral base-lattice map have not been discarded.
The sum after $m$ steps is $\sum_i B_i-(m-1)d_w$.
The theorem \texttt{nuIn\_mem\_Sram} places the indicated input in
that initial local lattice.  Finally \texttt{period\_ge}, using
\texttt{PairInt}, represents the actual terminal period over the
coefficient field and subtracts one further $d_w$.  Its elementary
valuation proof is the multiplication identity
$v_w(ay)=v_w(a)+v_w(y)$ for nonzero factors, with separate zero cases.
This proves \eqref{eq:SF08-pointwise}, including the ramification
factor and every stabilizer loss.
\end{proof}

\paragraph{Why the stripped hypotheses are legitimate.}
The packaged \texttt{Shape.rankA} condition would impose
$2\log_5(a^2)+2\le k-1$.  It fails for the candidate below, but
it is not used by the one-orbit proof just given.  Similarly the
packaged true-column theorem assumes $3\le\texttt{splitBound}$.
In its fixed inert-prime branch that assumption is used to obtain
$p\ge5$; here $p=863$ supplies this directly.
The actual standard-series theorem
\texttt{thetaExpansionRealPlaces\_stdSer\_of\_adm} has hypotheses
\texttt{HostAt}, $5\le p$, the residue-class and $h_P=0$ conditions,
admissible lattice pairs, standard bases, and the stated normalized
combination identity.  It has no \texttt{rankA} or
$3\le\texttt{splitBound}$ hypothesis.  The fixed-prime part of the
adapted-kernel proof can therefore be repeated using that theorem
and the actual identification of $E.\texttt{Th}$ with its standard
series.  This observation reduces an input gate; it does not
manufacture the required actual identification or integral bases.

\paragraph{The explicit two-step partitions, including a variable gap.}
Take $k=8$, $m=2$, $t=15$, $f=8$, so
$(N_0,N_1,N_2)=(8,24,40)$.  At every embedding take
\[
 u=\tfrac12(113,45,37,29,21,13,5,-1),\qquad
 (b_1,b_2)=(M,M+G),\qquad
 (e_1,e_2)=(M+18,M+G+18).
\]
The source formula \texttt{nuRow}, with its ordinary derivative
normalization, gives
\begin{align*}
 \nu_0={}&(M+59,M+26,M+23,M+20,M+17,M+14,M+11,M+9),\\
 \nu_1={}&(M+G+43,M+G+10,M+G+7,M+G+4,\\
          &\hspace{8mm}M+G+1,M+G-2,M+G-5,M+G-7,(G-16)^{16}).
\end{align*}
For integers $M\ge0$ and $G\ge16$ these are partitions, and their sizes are exactly
\[
 |\nu_0|=8M+179,\qquad |\nu_1|=8M+24G-205.
\]
These identities follow by substituting into the first eight rows
of \texttt{srcPar} and its sixteen added rows; in the latter rows
the two half-integer rank shifts cancel, leaving $G-16$.
Thus $\sum_{i,z}s_{i,z}=128M+192G-208$ and
$f(N_0N_1+N_1N_2)=9216$.
For $p=863$, SF-Q6-04(a) gives $d_w=0$ for the genuine denominator.
With $A=a^2$, \eqref{eq:SF08-pointwise} becomes
\begin{equation}\label{eq:SF08-explicit}
 R_G=\frac{128M+192G-208}{744768}
                      -9216(2\log_{863}a+1).
\end{equation}

If in addition the actual F/E/gain ties of SF-Q6-04(b) are supplied,
their full expression for this variable gap is
\[
 \operatorname{sg}_G
     =\frac{96M+48G+1728}{744768}-76800\log_{863}a.
\]
For $M\ge0$, $G\ge16$ and $\log_{863}a\ge1$ one has
\[
 R_G-\operatorname{sg}_G
 =\frac{32M+144G-1936}{744768}
                   +58368\log_{863}a-9216>0.
\]
Consequently the pointwise theorem implies this full weaker bound
when all its actual arithmetic ties hold.

For $B=10^{10}$, $M=\lceil e^B\rceil$,
$a=10^{1000}e^B$, and any $16\le G\le M$, the size hypotheses with
$A=a^2$ hold: $M\le2e^B$, $M\ge179$, and both partition sizes
are at most $32M$, whereas $a^2>32M$ and $a^2>863$.
For $G\ge36$, $\operatorname{sg}_G$ is at least the positive
expression of SF-Q6-04(c), whose explicit rational lower bound is
$22020927106278400000/3879$.
The old illustrative gap $G=36$ is sufficient for this pointwise
calculation; it is \emph{not} asserted to satisfy the separate
\texttt{NodeGap} hypotheses of theta-chain surjectivity.
A large gap, for example $G=10^9$, retains this local positivity.
For the exact \texttt{NodeGap} gate, $\operatorname{rank}X'=39$,
so its existential auxiliary dimension can be chosen as $N=40$.
The required bound is $15+2\cdot40^3+1=128016<10^9$.
Thus the large gap satisfies this particular gate.  This does not
supply the other geometric and analytic hypotheses of surjectivity.
The exact definition is
\texttt{Development/AbelianVarietyTheory/BlockTheory/}
\texttt{TwistedDiscreteSeries/ThetaSourceAction.lean}, lines 209--212,
with a run copy in the critical reviewer's \texttt{pu\_source\_cache}.

\paragraph{Remaining actual input gates.}
An actual global specialization still requires the relevelled host
and its actual column datum; the prescribed kernel and normalized
standard-series identification; integral frames supplying
\texttt{KForm}; the actual row spaces, twists, and integral
\texttt{PairInt}; and, to apply the complete DivInputs theorem,
its cutoff, row, gain, F/E, input-lattice and size-data ties.
A coefficient-field realization of a finite kernel does not prove
these theta-factorization or integral-pairing statements.
The constant-twist model is separately rejected by SF-Q6-07.
The theorem above therefore removes a real obstruction in the
\emph{packaged proof}, but is not a certified nonzero global period
or an unconditional global divisibility example.

\paragraph{Exact archive and run locators.}
All paths below are below \texttt{lean/src/} in the supplied archive;
run copies are below \texttt{lean\_selected/} or
\texttt{SIGN\_NEXT/work/singular\_fibres/arithmetic\_review/source\_cache/}.
\begin{itemize}
\item \texttt{Formalism/SharpRestriction/Divisibility/}
\texttt{ProductZerothFirstParts.lean}, lines 263--272:
\texttt{Loc1.Inp}, \texttt{HidaSetting}, \texttt{ThmOne}.
\item \texttt{JetPolynomialLogs/DworkPowerSeriesWeights.lean},
\texttt{thmOne\_of}, lines 235--266; and
\texttt{JetPolynomialLogs/LocalBoundElevenComparison.lean},
\texttt{split\_bound\_eleven\_one}, lines 64--68, below that
Divisibility directory.
\item \texttt{Formalism/SharpRestriction/Divisibility/ChainOperators/OrbitEntryBounds.lean},
\texttt{orbitsAddAtConjugate\_ent}, lines 23--24.
\item \texttt{Formalism/SharpRestriction/Divisibility/OrbitPlaces.lean},
\texttt{Ent}, lines 34--52; \texttt{KForm}, lines 76--79;
\texttt{compose}, lines 806--814; $n_P=1$, lines 1129--1137.
\item \texttt{Formalism/SharpRestriction/GrowthBounds/}
\texttt{WeilJacquetShalikaLaw.lean}, \texttt{PairInt}, lines 383--395;
\texttt{period\_ge}, lines 419--442;
\texttt{nuIn\_mem\_Sram}, beginning line 444.
\item \texttt{Development/AbelianVarietyTheory/Proofs/SatakeParameters/}
\texttt{B0StandardSeriesAdmissible.lean}, lines 26--39;
the full fixed-prime proof in
\texttt{FrobeniusConjugationConditionsAdaptedKernelDatum.lean};
and the packaged wrapper \texttt{AtOfTrueColumns.lean}.
\item \texttt{Formalism/LocalRingCardinality/DataDIVDefs.lean},
\texttt{srcPar} and \texttt{nuRow}, lines 60--73.
\end{itemize}
The displayed proof rechecks the complete valuation-composition,
terminal-pairing, and one-orbit wrapper proofs.  Deeper analytic
prerequisites of the imported theta expansion and one-letter estimate
are used as source theorems; no complete Lean build or independent
reproof of all transitive dependencies is claimed.
